\clearpage
\VerdierPhysicalPageAnchor{0250}%
\begingroup\small\setlength{\parskip}{0pt}\setlength{\emergencystretch}{1.5em}
\setlength{\abovedisplayskip}{4pt plus 2pt minus 2pt}%
\setlength{\belowdisplayskip}{4pt plus 2pt minus 2pt}%

\VerdierSemanticAnchor{ch4d-4-6-9-3-isomorphism}
\noindent Ce morphisme est un isomorphisme lorsque, par exemple, le
complexe $Y$ est un objet de $\operatorname{comp}^{+}(\mathcal A)$ et
lorsque le foncteur $F$ est stationnaire à droite, ou bien, sans hypothèse
sur $Y$, lorsque le foncteur $F$ est stationnaire (4.5.1).

\VerdierSemanticAnchor{ch4d-4-6-9-derived-functors}
Le double complexe $I^{\bullet\bullet}$ étant une résolution injective de
Cartan--Eilenberg de $Y$, le complexe
$H_1^r(I^{\bullet\bullet})$ est une résolution injective de l'objet
$H^r(Y)$ et, par suite, on a un isomorphisme :

\VerdierSemanticAnchor{ch4d-4-6-9-4}
\[
H_{\mathcal B}^{p}\left(
F^{q-r,\bullet}(H_1^r(I^{\bullet\bullet}))
\right)
\xrightarrow{\sim}
R^pF^{q-r}(H^r(Y)),
\tag{4.6.9.4}
\]
où $R^pF^{q-r}\colon\mathcal A\longrightarrow\mathcal B$ désigne le
$p$-ième dérivé droit (au sens de [1]) du foncteur $F^{q-r}$ restreint à
la catégorie $\mathcal A$. On a donc, lorsque le morphisme (4.6.9.3) est un
isomorphisme, un isomorphisme :

\VerdierSemanticAnchor{ch4d-4-6-9-5}
\[
(\mathrm{C.E.})_2^{p,q}\xrightarrow{\sim}
\prod_{r\in\mathbb Z}R^pF^{q-r}(H^r(Y)).
\tag{4.6.9.5}
\]

\VerdierSemanticAnchor{ch4d-4-6-10}
\medskip\noindent\textbf{4.6.10.}\quad
Étudions le cas où la catégorie $\mathcal B$ est la catégorie
$\mathbf{Ab}$ des groupes abéliens appartenant à un univers convenable et
où le foncteur :

\[
F\colon D(\mathcal A)\longrightarrow\mathbf{Ab}
\]
est le foncteur :

\VerdierSemanticAnchor{ch4d-4-6-10-ext0}
\[
Y\longmapsto F(Y)=\operatorname{Ext}^{0}(Z,Y),
\]
où $Z$ est un objet de $D(\mathcal A)$. Soient $Y$ un complexe de
$\mathcal A$ et
$\varepsilon\colon Y\longrightarrow I^{\bullet\bullet}$ une résolution
injective de Cartan--Eilenberg. Le morphisme (4.6.9.3) est alors un
isomorphisme (4.6.3) et, par suite, on a un isomorphisme :

\VerdierSemanticAnchor{ch4d-4-6-10-1}
\[
(\mathrm{C.E.})_2^{p,q}\xrightarrow{\sim}
\prod_{r\in\mathbb Z}
\operatorname{Ext}^{p}\bigl(H^{r-q}(Z),H^r(Y)\bigr).
\tag{4.6.10.1}
\]

\VerdierSemanticAnchor{ch4d-4-6-10-stationarity}
Par ailleurs, la suite spectrale de Cartan--Eilenberg relative à $F$ est
stationnaire lorsque $Y$ est un objet de
$\operatorname{comp}^{+}(\mathcal A)$ et $Z$ un objet de
$D^{-}(\mathcal A)$, ou bien lorsque la résolution injective de
Cartan--Eilenberg est de longueur finie. De plus, sous ces hypothèses,
l'aboutissement de cette suite spectrale est, en vertu de (4.6.8) :

\VerdierSemanticAnchor{ch4d-4-6-10-abutment}
\[
\prod_{r\in\mathbb Z}
\operatorname{Ext}^{p}\bigl(H^{r-q}(Z),H^r(Y)\bigr)
\Longrightarrow
\operatorname{Ext}^{p+q}(Z,Y).
\]
\endgroup
